``Big data'' refers to data that are more voluminous, but often also more varied, unstructured, and dynamic than data traditionally used in communication research. In Communication Science and the Social Sciences more broadly, this includes research using internet-based data sources such as social-media content, digital news archives, platform data, and public comments. This area of research is commonly described as \emph{Computational Social Science} \citep{Lazer2009} or, when focused on communication processes and content, \emph{Computational Communication Science} \citep{Shah2015}.

This course introduces the concepts, opportunities, and challenges of computational research with digital communication data. Computational approaches can make it possible to analyse large or complex collections of texts and other digital traces; they complement rather than replace established communication-science methods such as theory-driven research design, manual content analysis, and statistical inference. Throughout the course, we emphasise that automated outputs require substantive interpretation and empirical validation.

Participants learn practical strategies for acquiring, organising, preprocessing, analysing, and documenting digital data in communication contexts. The course focuses on (a) data acquisition from files and APIs, data wrangling, exploratory analysis, and reproducible workflows; and (b) automated content analysis, including text preprocessing, dictionary and Bag-of-Words approaches, unsupervised text analysis, embeddings, topic modelling, and supervised machine-learning classification. Students learn to inspect and evaluate automated outputs, compare them with appropriate benchmarks, and reflect on validity, ethics, reproducibility, and the responsible use of AI tools.

The course uses Python as its main programming environment. Students are expected to be willing to learn to write, read, adapt, and explain Python code. Previous programming experience is helpful but not required. Students without prior experience are encouraged to complete one of the many free online introductions to Python before the course begins.
